% !TEX root = ../main.tex
% anything appart from comments in this file will appear in the Research Section of the text.
% do your main writing here.
% use \autocite for bibliography and \ref to cite internal figures and tables

% Detailed research into the style/form in question, referencing a wide range of relevant source material. Clear evidence of critical reflection and analysis of sources that informs both aesthetic and functional aspects to the work.

The main inspiration for this was S1E6 [citation] of \textit{Arcane} (2021), "When These Walls Come Tumbling Down". In the episode's climax, the protagonists Vi and Caitlyn flee from antagonist Silco's goons and Silco's plans come crashing down around him (figuratively and literally). This is accompanied not by frenetic action music, but by Woodkid's "Guns for Hire", a sombre, yet hopeful, downtempo track, consisting primarily of voice, guitar, and strings.

I think this choice of music is so effective for the juxtaposition of the action on the screen with the tempo and feel of the song. As Vi's long-estranged sister lights a flare for Vi to find her, the music swells into a sombre, yet hopeful chorus [insert time stamp], and even as Silco viciously and repeatedly kicks one of his subordinates bloodier and bloodier, the same song continues.

This, of course, is nothing new in film and TV, and this kind of juxtaposition between the on-screen action and the soundtrack is a tried and tested formula. It's cliche, but that means it's effective.

Now, a TV-show might not be the most obvious choice of reference for a video game cutscene, but the same kind of juxtaposition can work just as well here. In fact, you could take "Guns for Hire" exactly as it is presented on the soundtrack album and play it over the cutscene, and it works quite well as-is.

But still, using a song to score a video game opening cutscene is --- as far as I know --- unusual. However, using songs in video games has become more common as the hardware limitations of game consoles have lessened. Darren Korb has had success doing just that with, among others, \textit{Bastion} (2011) [check this and get a reference], \textit{Hades} (2020) and \textit{Hades II} (2025) [check this too]. Additionally, Lorien Testard's score for \textit{Clair Obscur: Expedition 33} (2025) uses sung vocals extensively throughout the game to great effect.

\subsection*{"Guns for Hire" highlights / analysis}

The track uses a lot of clever/interesting harmonic shifts that I have taken the liberty of either slightly modifying or stealing outright.

It's ostensibly in B-flat minor, but the occasional major second chord (C major in this case) is featured prominently, which gives a sense of foreboding, probably due to the raised fourth scale degree being a tritone above the root.  [get example of where this is in the track and where I've used the same one]

Another very cool trick is the sudden tritone jumps ["you play with your blocks"] which suddenly sheds the darkness, especially when it jumps to a major chord and runs down a string of major chords [this needs closer checking]

Finally, the vocals landing on the major seventh of the major sixth chord  [first verse, chorus] is a repeated motif that paints the track in a bittersweet light. The minor tonic chord to the major sixth chord really has a pretty hopeful sound, but with the vocals landing hard on the major seventh, it adds a delicate piece of instability to the whole thing.
